% ------------------------------------------------------------------------------
% Formatvorlage für Masterarbeiten an der Ohm-Hochschule N�rnberg
% ------------------------------------------------------------------------------
% Dokumentenkopf ---------------------------------------------------------------
%   Diese Vorlage basiert auf "scrreprt" aus dem koma-script.
% ------------------------------------------------------------------------------
\documentclass[
    12pt, % Schriftgröße
    DIV=10,
    % ngerman, % fürUmlaute, Silbentrennung etc.
    a4paper, % Papierformat
    oneside, % einseitiges Dokument
    titlepage, % es wird eine Titelseite verwendet
   	%parskip=half, % Abstand zwischen Absätzen (halbe Zeile)
   	%parskip=full-,
   	parskip=false,
    %headings=normal, % Größe der Überschriften verkleinern\l
    toc=chapterentrydotfill,
    listof=totoc, % Verzeichnisse im Inhaltsverzeichnis aufführen
    bibliography=totoc, % Literaturverzeichnis im Inhaltsverzeichnis aufführen
    index=totoc, % Index im Inhaltsverzeichnis aufführen
    numbers=noenddot, %Gliederung wird ohne Punkt nach Gliederungszahl dargestellt
    captions=tableheading, % Beschriftung von Tabellen unterhalb ausgeben
    final % Status des Dokuments (final/draft) 
]{scrreprt}

% Hier sind Metainformationen drinnen
% Titel, Betreuer, Studiengang etc
% Meta-Informationen -----------------------------------------------------------
%   Definition von globalen Parametern, die im gesamten Dokument verwendet
%   werden können (z.B auf dem Deckblatt).
%
% ------------------------------------------------------------------------------
\newcommand{\titelen}{Development and Evaluation of a Lifecycle Management Interface for Improving Stability and Integration of Mobile Robots based on ROS2}
% \newcommand{\titelde}{Systemintegration für mobile Robotik}

% \newcommand{\titel}{Systemintegration for mobile Robotics}

%\newcommand{\titel}{Developing a Lifecycle Management Interface for improving System-stability and System-integration for mobile ROS2 robotic systems using managed lifecycle nodes, docker and siemens industrial edge.}

\newcommand{\titel}{Development and Evaluation of a Lifecycle Management Interface for Improving Stability and Integration of Mobile Robots based on ROS2}

\newcommand{\art}{Masterthesis}
\newcommand{\autor}{Aayush Yadav}
\newcommand{\hochschule}{Technische Hochschule Nürnberg Georg Simon Ohm}
\newcommand{\fakultaet}{Elektrotechnik Feinwerktechnik Informationstechnik}
\newcommand{\studiengang}{Elektronische und Mechatronische Systeme}
% \newcommand{\studienbereich}{Automatisierungstechnik}
\newcommand{\matrikelnr}{3549452}
\newcommand{\erstgutachter}{Prof. Dr. Stefan May}
\newcommand{\zweitgutachter}{}
\newcommand{\jahr}{2022}
\newcommand{\abgabedatum}{20.07.2022}
\newcommand{\semester}{Summer Semester 2022}
\newcommand{\firma}{SIEMENS}
\newcommand{\firmenbetreuerA}{Michael Fiedler}
\newcommand{\firmenbetreuerB}{Florian Gramss}
\newcommand{\ort}{Nürnberg}
\newcommand{\logoOhm}{LogoOhmHochschule.jpg}
\newcommand{\logoSIEMENS}{SIEMENSLogo.jpg}



% Hier sind alle Pakete drinnen
% Seitenränder
\usepackage{geometry}

% Zeilenabstand
\usepackage{setspace}

% Koma Script - Seitenlayout
\usepackage[
    automark, 		% Kapitelangaben in Kopfzeile automatisch erstellen
	headsepline, 	% Trennlinie unter Kopfzeile
	ilines 			% Trennlinie linksbündig ausrichten
]
{scrlayer-scrpage}

% Koma Script - Inhaltsverzeichnis
\usepackage{tocbasic}

% Silbentrennung
% \providehyphenmins{ngerman}{44}
\providehyphenmins{english}{44}

% Sprachpaket
% \usepackage[ngerman]{babel}
\usepackage[english]{babel}

% Zum einbinden anderer PDF Dokumente
\usepackage{pdfpages}

% Generiert sinnlosen Text
\usepackage{lipsum}

% Hübschere Tabellen
\usepackage{tabularx}
\usepackage{longtable}

% Bilder und Tabellen an exakt dieser Position ( [H] ):
\usepackage{float} 		

% Sub-figures 
\usepackage{caption}
\usepackage{subcaption}

% DVD-Verzeichnis darstellen
\usepackage{dirtree}

% Listings darstellen
\usepackage{listings}
\usepackage{lstautogobble}

% Grafiken einbinden
\usepackage{graphicx}
% Biber (Biblography)
\usepackage[
	style=numeric, % Loads the bibliography and the citation style 
	% bibstyle=alphabetic, % load a bibliography style
	% citestyle=alphabetic, % load a citatio style
	sorting=none,
	natbib=true, % define natbib compatible cite commands
	%%--- Backend --- --- ---
	backend=biber,   % (bibtex, biber)
	bibwarn=true,     %
	texencoding=auto, % auto-detect the input encoding
]{biblatex}  
% Other options:
%  style=numeric, % 
%  style=numeric-comp,    % [1-3, 7, 8]
%  style=numeric-verb,    % [2]; [5]; [6]
%  style=alphabetic,      % [Doe92; Doe95; Jon98]
%  style=alphabetic-verb, % [Doe92]; [Doe95]; [Jon98]
%  style=authoryear,      % Doe 1995a; Doe 1995b; Jones 1998
%  style=authoryear-comp, % Doe 1992, 1995a,b; Jones 1998
%  style=authoryear-ibid,
%  style=authoryear-icomp,
%  style=authortitle,
%  style=authortitle-comp,
%  style=authortitle-ibid,
%  style=authortitle-icomp,
%  style=authortitle-terse,
%  style=authortitle-tcomp,
%  style=authortitle-ticomp,

% UML und Sequenzdiagramme
\usepackage{pgf-umlsd}


% Für schöne, selbstgezeichnete graphiken: tikz
\usepackage{tikz}
\usetikzlibrary{patterns}
\usetikzlibrary{positioning}
\usetikzlibrary{shapes,arrows}
\usetikzlibrary{arrows.meta}
\usetikzlibrary{shapes.geometric, arrows}
\usetikzlibrary{calc}
\usetikzlibrary{decorations.pathreplacing}


\usepackage{fontawesome}






% Hyperref für Verlinkungen 
% ACHTUNG: Als letztes Paket ! (Mit ausnahme von Glossaries)
% PDF-Optionen -----------------------------------------------------------------
\usepackage[
	bookmarks,
	bookmarksopen=true,
	colorlinks=true,
	% diese Farbdefinitionen zeichnen Links im PDF farblich aus
	% linkcolor=red, % einfache interne Verknüfungen
	% anchorcolor=black,% Ankertext
	% citecolor=blue, % Verweise auf Literaturverzeichniseinträge im Text
	% filecolor=magenta, % Verkrüfungen, die lokale Dateien öffnen
	% menucolor=red, % Acrobat-Menüpunkte
	% urlcolor=cyan, 
	% diese Farbdefinitionen sollten für den Druck verwendet werden (alles schwarz)
	linkcolor=black, % einfache interne Verknüpfungen
	anchorcolor=black, % Ankertext
	citecolor=black, % Verweise auf Literaturverzeichniseinträge im Text
	filecolor=black, % Verknüfungen, die lokale Dateien öffnen
	menucolor=black, % Acrobat-Menüpunkte
	urlcolor=black, 
	%pagebackref,
	plainpages=false, % zur korrekten Erstellung der Bookmarks
	pdfpagelabels, % zur korrekten Erstellung der Bookmarks
	hypertexnames=false, % zur korrekten Erstellung der Bookmarks
	%linktocpage % Seitenzahlen anstatt Text im Inhaltsverzeichnis verlinken
	linktoc=all
]{hyperref}

\hypersetup{
	pdftitle={\titel},
	pdfauthor={\autor},
	pdfcreator={\autor},
	pdfsubject={\titel},
	pdfkeywords={\titel},
}

% Glossar und Abkürzungsverzeichnis, laut doku nach Hyperref
\usepackage[
	acronym, 
	nonumberlist, 
	toc,
	numberline,
	nopostdot,
	nogroupskip
]{glossaries}


% In Seitenstil wird das Layout der Seiten festgelegt
% Sachen wie Seitenränder, Zeilenabstand, Layout Inhaltsverzeichnis
% Seitenränder
\geometry{
	paper=a4paper,
	right=20mm,		
	top=20mm,
	left=30mm,
	bottom=30mm
}

% Zeilenabstand 1,5 Zeilen
\singlespacing
%\spacing{1.21}
%\onehalfspacing
%\doublespacing


% Minimal mehr platz nach einem punkt. <-
\frenchspacing

% Schusterjungen und Hurenkinder vermeiden
\clubpenalty = 10000
\widowpenalty = 10000 
\displaywidowpenalty = 10000


% Style der listings
\definecolor{RoyalBlue}{cmyk}{1, 0.50, 0, 0}
\definecolor{codegreen}{rgb}{0,0.6,0}
\definecolor{codegray}{rgb}{0.5,0.5,0.5}
\definecolor{codepurple}{rgb}{0.58,0,0.82}
\definecolor{backcolour}{rgb}{0.95,0.95,0.92}
\lstset{
	basicstyle=\ttfamily,
	numberstyle=\tiny\color{gray},
	keywordstyle=\color{RoyalBlue},
	commentstyle=\color{green}\ttfamily,
	stringstyle=\color{mauve},
	autogobble=true,
	breaklines=true,
	breakatwhitespace=true,
	captionpos = b,
	frame=single, % box um das listing
	numbers=left, %nummerierung um das Listing
	% Nummerierung innerhalb der box
	framexleftmargin=5.5mm,
	xleftmargin=19pt,
	xrightmargin=3.4pt,
} 


%Ligaturen verhindern:
%\DisableLigatures{}
%Was sind ligaturen ?:
%ff und fi ragen ineinander, und auch andere wilde Buchstabenkombinationen


% Kopf und Fußzeile
\pagestyle{scrheadings}
% Kopf- und Fußzeile auch auf Kapitelanfangsseiten
\renewcommand*{\chapterpagestyle}{scrheadings}
% Schriftform der Kopfzeile
\renewcommand{\headfont}{\normalfont}

% Höhe der Kopfzeile
\setlength{\headheight}{14mm}
% Kopfzeile
\ihead{\headmark}
\chead{}
%\ohead{\includegraphics[scale=0.15, trim={2.4cm 24.5cm 2.6cm 2.6cm}, clip]{Bilder/SIEMENS-Logo.pdf}}
\ohead{\includegraphics[scale=0.08]{Bilder/Siemens_Logo.pdf}}
%\ohead{\art} % Masterarbeit statt SIEMENS in der Kopfzeile

% Fußzeile
% Leer, da sich das Layout ändert
\cfoot{}
\ifoot{}
\ofoot{}

% Layout Inhaltsverzeichnis
\setcounter{tocdepth}{2}  % Inhaltsverzeichnis nur bis 2 Unterelemente auflisten
\setcounter{secnumdepth}{3} % Numerierungstiefe für Inhaltsverzeichnis

\DeclareTOCStyleEntries[
% numwidth=1em,					% Abstand Zahl - Überschrift
rightindent=4em,				% Rand nach Links. Relevant für Überschriften mit Zeilenumbruch im Inhaltsverzeichnis
pagenumberbox=\pagenumberbox
]
{tocline}{chapter, section, subsection, subsubsection, paragraph, subparagraph, figure, table}
\newcommand*\pagenumberbox[1]{\mbox{\hspace{2pt}#1}} 

%Höhe von Kapiteln im Text					2.3
%\renewcommand*{\chapterheadstartvskip}{\vspace*{1.0\baselineskip}}% Abstand einstellen
\renewcommand*{\chapterheadstartvskip}{\vspace*{1\baselineskip}}% Abstand einstellen

% Eigenes Design für Abkürzungsverzeichnis
\newglossarystyle{mylong}{%
	\renewenvironment{theglossary}%
	{\begin{longtable}[l]{llp{\glsdescwidth}p{\glspagelistwidth}}}%
		{\end{longtable}}%
	\renewcommand*{\glossaryheader}{}%
	\renewcommand*{\glsgroupheading}[1]{}%
	\renewcommand{\glossentry}[2]{%
		\textsf{\textbf{\glsentryitem{##1}\glstarget{##1}{\glossentryname{##1}}}} &		\glossentrysymbol{##1} &		\glossentrydesc{##1} &		##2\tabularnewline
	}%
	\renewcommand{\subglossentry}[3]{%
		&
		\glssubentryitem{##2}%
		\glossentrysymbol{##2} &
		\glstarget{##2}{\strut}\glossentrydesc{##2} & ##3\tabularnewline
	}%
	\renewcommand*{\glsgroupskip}{%
		\ifglsnogroupskip\else & & &\tabularnewline\fi}%
}

% Abstand zwischen Bild und Bildunterschrift
% \setlength\abovecaptionskip{-20pt}
%Vektorbasierte Standardschrift
\renewcommand{\familydefault}{\sfdefault} 	% %Defaultschrift ohne Serifen (modernen)
\usepackage{sansmath}						% %Matheschrift ohne Serifen
\usepackage{microtype}						% %Verbesserte Textverteilung in Zeile

\lstdefinelanguage{JavaScript}{
  keywords={typeof, new, true, false, catch, function, return, null, catch, switch, var, if, in, while, do, else, case, break},
  keywordstyle=\color{blue}\bfseries,
  ndkeywords={class, export, boolean, throw, implements, import, this, import},
  ndkeywordstyle=\color{codepurple}\bfseries,
  identifierstyle=\color{black},
  sensitive=false,
  comment=[l]{//},
  morecomment=[s]{/*}{*/},
  commentstyle=\color{purple}\ttfamily,
  stringstyle=\color{red}\ttfamily,
  morestring=[b]',
  morestring=[b]"
}


% Biblographie hinzufügen
\addbibresource{Quellen/Quellen_Online.bib}
\addbibresource{Quellen/Quellen_Buecher.bib}
\addbibresource{Quellen/Quellen_Verschiedenes.bib}
% Glossar erstellen
% Index erstellen
\makeindex
% Datei einbinden
\loadglsentries{Inhalt/00_Glossar.tex}
% Glossar generieren
\makeglossaries

% Da Von Römisch - Arabisch - Römisch gewechselt wird, wird ein Speicher benötigt
\newcounter{SeitenzahlSpeicher}



\begin{document}
	 % Leider muss laut vorgabe zwischen zwei Layouts gewechselt werden...
	 % Daher auch hier einzelne Befehle die nicht einheitlich für das gesamte Dokument sind

	
	% Sprache festlegen
	\selectlanguage{english}
	% \englishcontentsname{Table of Contents}

	% Leerseite und Deckblatt
	%Leerseite
\newpage
\thispagestyle{empty}
\section*{}
\newpage

%Deckblatt
\thispagestyle{empty}

\newgeometry{left=25mm, right=15mm, top=15mm, bottom=5mm}
\begin{titlepage}
	% Textteil
	\begin{center}
		
		\includegraphics[width=\textwidth]{Bilder/TH_Logo.pdf}\\[0.2cm]
		\textsc{\LARGE Fakultät efi}\\[1ex]
		\Large{Studiengang \studiengang}\\
		% \Large{Vertiefungsrichtung \studienbereich}\\
		
		
		% Titel
		\vspace*{40mm}
		\hrule
		\vspace*{5mm}
		\LARGE{\textbf{\titelen}}
		\vspace*{3mm}
		\hrule
		
		% Name
		\vspace*{15mm}
		\Large{{\art} von}\\
		\LARGE{\textsc{\autor}}\\
		\Large{Matrikel-Nr.: \matrikelnr}\\
		
		
		% Zusätzliche Infos
		\vspace*{20mm}		
		\Large{\semester}\\		
		\Large{Abgabedatum: \abgabedatum }\\
		
		\vfill	
		\large{Betreuer:\hfill   \mbox{}}\\
		\large{\erstgutachter \hfill \firmenbetreuerA \mbox{}}\\	
		%\large{Zweitgutachter: \zweitgutachter \hfill \firmenbetreuerB \mbox{}}\\
		\large{Technische Hochschule Nürnberg \hfill\firma}\\[2ex]
		
		% Schlagworte
		\large{Schlagworte: Docker, Industrial Edge, ROS2\hfill   \mbox{}}
		
	\end{center}
	
	
\end{titlepage}
\restoregeometry

	
	
	% Hier mit Römisch Seitenzählen anfangen
	\pagenumbering{Roman}
	\setcounter{page}{2}
	
	% Prüfungsrechtliche Erklärung
	%\includepdf[pages={1}]{SB_0500_Ausgefuellt.pdf} %Erklärung zur veröffentlichung
	
	% Hier Römische Seitenzahlen auch angeben
	\ofoot{\pagemark}
	
	% Abstract
	% Inhaltsverzeichnis
 	\tableofcontents
	
	% Glossar
	% Alle Einträge ausgeben
	\glsaddall
	% Glossar darstellen
	\printglossary
	% \label{sec:Glossar}
	
	% Abkürzungsverzeichnis
	\printglossary[type=\acronymtype, title=Abkürzungsverzeichnis, style=mylong]

	% Aufräumen und Seitenanzahl speichern
   	\clearpage
   	\setcounter{SeitenzahlSpeicher}{\value{page}}
   	\label{Ende_Vorgeplaenkel}
   	
	% Inhalt
	% Hier wieder umstellen auf arabisches Seitenzählen
	\pagenumbering{arabic}
	\ofoot*{{\pagemark} von \pageref{Ende}}
	
	\chapter{Einleitung}
\label{Einleitung}

	\section{Motivation}
	\label{Einleitung:Motivation}
	
	\section{Aufgabenstellung}
	\label{Einleitung:Aufgabenstellung}
	asdf fdsa\\
	Nach \cite[S. 169]{Krypczyk2018} müssen folgende Fragen nach der Analysephase des Softwareentwicklungsprozess beantwortbar sein:
	\begin{itemize}
		\item Was sind die entscheidenden Anforderungen an die zu erstellende Software?
		\item Welches Problem soll mithilfe des Anwendungssystems gelöst werden?
		\item Welche Wünsche haben Ihre Kunden und Nutzer an das System?
	\end{itemize}
	
	
	% \chapter{Stand der Technik}
\label{StandDerTechnik}
	\section{Bestehende Systeme}
	\label{StandDerTechnik:Sub1}
	
	\section{Zusammenführen von Programmen jetzt}
	\label{StandDerTechnik:Sub2}
	Diese Masterarbeit 
	Das System \cite{Laura}
	
	% \chapter{Grundlagen}
\label{Grundlagen}
Im Rahmen dieser Masterarbeit werden mehrere Softwarepakete und bereits bestehende Programme verwendet. Nachfolgend werden die wichtigsten Programme erläutert.
	\section{Docker}
	\label{Grundlagen:Docker}
		Docker bietet die Möglichkeit Anwendungen auf einem System zu isolieren und Strukturieren.
		Hiermit lassen sich einzelne Programme, oder zum Zeitpunkt der Containererstellung festgelegte Prozesse kapseln, in einen Container überführen und somit eine hochgradig agile Umgebung schaffen. \cite{dockerGetStarted}.
		In Abbildung \ref{fig:Grundlagen:Docker:DockerVsVM} wird der Unterschied zwischen einer \gls{VM} und der Docker Umgebung dargestellt.
		Es ist ersichtlich, dass innerhalb eines Docker Containers nicht von Grund auf das Betriebssystem ausgeführt werden muss. 
		Somit lassen sich die zur Verfügung stehenden Ressourcen des Host Computers effizienter bereitstellen.
		\begin{figure}[H]
			\centering
			\includegraphics[width=\textwidth]{"Bilder/DockerVsVM.png"}
			\caption{Vergleich zwischen Docker und Virtuellen Maschinen. \cite{dockerVsVM}}
			\label{fig:Grundlagen:Docker:DockerVsVM}					
		\end{figure}
		Um tiefer ins Detail zu gehen, ist die Erläuterung der Begriffe Image, Container, Dockerfile und Docker-Compose notwendig.
		\paragraph{Image}
			Ein Docker-Image ist ein Speicherabbild eines Containers. 
			Das Image besteht aus mehreren Layern, die schreibgeschützt sind.
			Images können manuell erstellt werden, oder mit Verwendung eines dockerfiles automatisiert kompilliert werden.
			
		\paragraph{Container}
			Ein Container ist die aktive ausgeführte Instanz eines Images.
			Es läuft immer ein Prozess innerhalb des Containers.
			Somit wird ein Container beendet, sobald er kein Programm ausführt oder das Programm mit seinem Auftrag fertig ist.
			
		\paragraph{Dockerfile}
			Ein Dockerfile ist eine Textdatei mit Anweisungen zum erstellen eines Containers.
			Jeder Schritt erstellt einen neuen Layer, der dem Image hinzugefügt wird.
			Mit einem Dockerfile lässt sich das erstellen von Images automatisieren.
		
		\paragraph{Docker-Compose}
			Mit Docker-Compose lassen sich die Startparameter eines Containers in einer .yaml -Datei festlegen.
			Es lassen sich mehrere Docker-Container gleichzeitig starten, sowie umfangreiche Einstellungen bei dem Start des Containers treffen.
		
		\subsection{Netzwerke}
		\label{Grundlagen:Docker:Netzwerke}
			Docker bietet auch die Netzwerkkapselung für die Container an.
			Somit lässt sich das Netzwerk eines Containers von dem Rest des Systems oder anderen Containern trennen.
			Es können auch eigene Netzwerke definiert werden, um mehrere Container miteinander Kommunizieren zu lassen.
			Sie spannen so ein eigenes Netzwerk auf, was vergleichbar mit einem lokalen Netzwerk ist.
			Zwei Netzwerktypen die Docker anbietet, werden nachfolgend erläutert.
			\paragraph{Bridged}
				Das Bridged Netzwerk kapselt das Netzwerk des Containers von dem des Host Systems, bietet aber noch eine Netzwerk-Brücke, damit der Container Teilnehmer in höherliegenden Netzen erreichen kann.
				Es lassen sich Ports des Host-Computers an einen Container weiterleiten, um so eine Erreichbarkeit des Containers aus einem höherliegendem Netz zu ermöglichen.
				
			\paragraph{Host}
				Die Netzwerkeinstellung Host verzichtet auf eine Kapselung der Netze.
				Hier verwendet der Container die IP Adresse des Host-Computers und ist entsprechend teil des Host-Netzes.
				Für weitere Teilnehmer des Netzwerkes, ist der Host-PC des Docker-Containers und der Container selbst, dasselbe Gerät.
				Eine Portweiterleitung zum ermöglichen von Diensten ist hier nicht notwendig.
				
			\subsection{Datenspeicherung}
			\label{Grundlagen:Docker:Datenspeicherung}
				Nach beenden des Containers werden die Daten, die während dem ausführen des Containers angefallen sind, gelöscht.
				Um eine Persistente Speicherung, zu ermöglichen, bietet Docker mehrere Möglichkeiten an, um über die Laufzeit des Containers hinaus Daten zu speichern, oder auch um auf Daten innerhalb des Containers zuzugreifen. 
				Nachfolgend werden die wichtigsten kurz vorgestellt.
				
				\begin{figure}[H]
					\centering
					% \input{"Bilder/tikz/Grundlagen/docker-types-of-mounts.tex"}
					\includegraphics[width=0.70\textwidth]{"Bilder/DockerSpeichertypenpng.png"}
					\caption{Verschiedene Datenspeichermöglichkeiten. \cite{dockerStorage}}
					\label{fig:Grundlagen:Docker:Datenspeicherung}					
				\end{figure}
										
				\paragraph{Volumes}
					Volumes werden auf dem Dateisystem des Hosts gespeichert (In Linux unter \textit{/var/lib/docker/volumes/}).
					Ein Volume kann in mehrere Container eingebunden werden, dadurch ist eine containerübergreifende Dateiablage möglich.
					Volumes werden nicht automatisch gelöscht, sobald alle Container gestoppt sind, welches das Volume eingebunden hatten.
					Sie eignen sich am besten, um Daten Persistent zu speichern. \cite{dockerStorage}
				
				\paragraph{Bind mounts}
					Bind mounts können überall auf dem Hostsystem gespeichert sein.
					So können wichtige Systempfade (zum Beispiel \textit{/dev/}) dem Container zugänglich gemacht werden. \cite{dockerStorage}
				
				\paragraph{tmpfs mounts}
					Hier werden die Dateien nicht auf dem Dateisystem des Hosts gespeichert, sondern auf dem Arbeitsspeicher des System.
					Dadurch entsteht eine nicht persistente Datenspeicherung, um geheime Informationen zwischen Containern auszutauschen. \cite{dockerStorage}
					
				
	\section{Industrial Edge}
	\label{Grundlagen:IndustrialEdge}
		Industrial Edge ist eine Entwicklung von Siemens, die ermöglicht, im Prozess anfallende Daten bereits auf dem Gerät zu analysieren.
		Es wird lokales Engineering mit Cloud Engineering kombiniert. \cite{siemensIEM_gettingStarted}
		Dadurch werden unnötige Datentransporte zwischen dem Endgerät und Server vermieden und erst bereits aufbereitete Daten zum Server geschickt.
		Die Rechenlast fließt so vom Server in Richtung des Endgerätes.
		Hierfür werden Applikationen (Apps) bereitgestellt, die einen festen Funktionsumfang enthalten und mit der Industrial Edge Infrastruktur auf das jeweilige Endgerät (Edge Device) gespielt werden können.
		Diese Apps basieren auf der Docker Technologie.
		Industiral Edge ermöglicht die zentrale Verwaltung von
		In der nachfolgenden Abbildung \ref{fig:Grundlagen:IndustrialEdge:Ueberblick} ist eine typische Infrastruktur von dem obersten Server, bis zu den Endgeräten dargestellt, die im Nachfolgenden näher erläutert wird.
		\begin{figure}[h]
			\centering
			%\input{"Bilder/tikz/Grundlagen/industrial-edge-overwiev.tex"}
			\includegraphics[width=0.70\textwidth]{"Bilder/Edge_uebersicht.jpg"}
			\caption{Überblick über Industrial Edge. \cite{siemensIEM_gettingStarted}}
			\label{fig:Grundlagen:IndustrialEdge:Ueberblick}					
		\end{figure}
	
		%- Management tool für Docker Container auf Geräten (Device)\\
		%- Graphen die das IE-Hub, IEM, Docker Device erklärt\\
			
		\paragraph{Industrial Edge Hub}
			Das \gls{IEH} ist die oberste, von Siemens bereitgestellte Serverebene.
			Hier bietet Siemens eigene entwickelte Applikationen an, die gegen eine Lizenzgebühr auf das \gls{IEM} geladen werden können.
			Von hier aus lassen sich ebenfalls benötigte Softwarepakete und Dokumentation herunterladen, die das Entwickeln von eigenen Apps und das Aufsetzen und betreiben von einem eigenen \gls{IEM} ermöglicht. \cite{siemensIEM_gettingStarted}
			
		\paragraph{Industrial Edge Management System}
			Das \gls{IEM} ist ein Server, welcher in einem Lokalen Netz selbstständig betrieben werden kann.
			Der Anwender hat hiermit die Möglichkeit, einen Server zu erstellen, damit z.B. sensible Daten, die zwischen dem Edge Device und dem \gls{IEM} ausgetauscht werden, nicht über das Internet transportiert werden müssen.
			Zudem lassen sich die Endgeräte über diesen Server verwalten, es können Apps sowie Softwareupdates aufgespielt werden, sowie lassen sich weitere Analysen durchführen.
			Eine selbst entwickelte App lässt sich auf das \gls{IEM} aufspielen, welches Diese wiederum an die Edge Devices verteilen kann.
			\cite{siemensIEM_gettingStarted}
			
		\paragraph{Industrial Edge App}
			Eine Edge App dient zur intelligenten Verarbeitung von Automatisierungsaufgaben. \cite{siemensIEM_gettingStarted}
			Neben den von Siemens im \gls{IEH} angebotenen Apps lassen sich auch eigene Applikationen entwickeln.
			Diese und eigene entwickelte Apps sind Docker basierte Images, welche auf dem \gls{IED} ausgeführt werden.
			Als oberste Beschreibungsebene wird eine docker-compose Datei angelegt, in welcher beschrieben wird, wie die docker Images zu starten sind, sowie welche Netzwerkkonfiguartion (Kapitel \ref{Grundlagen:Docker:Netzwerke}) und welche Art der Datenspeicherung (Kapitel \ref{Grundlagen:Docker:Datenspeicherung}) verwendet wird.
	
		\paragraph{Edge Device}
			Ein \gls{IED} führt die einzelnen Edge Apps aus. Sie können Automatisierungsdaten lokal speichern und nach Bedarf abrufen. 
			Ebenfalls lassen sich diese Daten an eine Cloud-Infrastruktur übermitteln. Ein \gls{IED} muss bei einem \gls{IEM} aktiviert werden.
			\cite{siemensIEM_gettingStarted}
			
		\paragraph{Industrial Edge Publisher}
			Der Industrial Edge Publishers ist ein Tool, mit dem Docker Images in Edge Apps überführt und in das \gls{IEM} hochgeladen werden.
			Es lassen sich ebenfalls bereits erstellte Apps verwalten, verändern oder löschen.
			\cite{siemensIE_App}

	
	\section{Simatic Automation Tool}
	\label{Grundlagen:SAT}
		Mit dem \gls{SAT} kann eine SPS konfiguriert werden.
		Dieses Tool steht für Windows sowie Linux mit deb oder rpm Paketverwaltung zur Verfügung. \cite{siemensSAT}
		Mit der Linux Version lässt sich die IP Adresse der SPS anpassen, ein bereits kompiliertes Programm downloaden, eine CPU in RUN oder STOP versetzen, eine CPU auf Werkseinstellung zurücksetzen, sowie die Firmware der SPS updaten\cite{siemensSAT_Online}. 
		Die Linux Version von \gls{SAT} stellt unter anderem folgende Befehle zur Verfügung:
		\begin{lstlisting}[caption={Beispielbefehle für SAT}, label={code:SATBeispiel}]
			transfer_sat -f /path/to/OMS.zip -i 123.231.123.3
			transfer_sat -c /path/to/OMS.zip
		\end{lstlisting}
		In Zeile eins des Listings \ref{code:SATBeispiel} wird das Kompilat welches in \lstinline|/path/to/OMS.zip| auf die SPS mit der IP Adresse \lstinline|123.231.123.3| geladen. 
		Voraussetzung ist, dass diese IP Adresse erreichbar ist.
		Mit dem Befehl in Zeile zwei lässt sich die angegebene Datei überprüfen.
		Anhand des Rückgabewerts lässt sich bestimmen, ob die angegebene Operation erfolgreich verlief.
		Hierbei wurde der für Linux etablierte Standard, dass ein Rückgabewert 0 einen Erfolg des Programms bestätigt, eingehalten.
		Weitere Fehlercodes sind unter anderem 100, der auf einen nicht erreichbaren Host hinweist, oder 103, welcher auf einen nicht erfolgreichen Download hinweist. \cite{siemensSAT}
					

	\section{ROS 2}
	\label{Grundlagen:ROS2}
		Ein asdfasdasd
		\cite{ros2Basic}
		\gls{ROS2}
	
	
	
		\subsection{Nachrichtentypen}
		\label{Grundlagen:ROS2:Nachrichtentypen}
			- Message (topics)\\
			- Service\\
			- Action
			
		\subsection{ROS 2-web Bridge}
		\label{Grundlagen:ROS2:2WebBridge}
			- Stellt einen websocket bereit, der mittels JSON requests verarbeitet. 
			- Basiert auf RclNodeJS, welches eine brücke zwischen javascript und ros2 bildet\\
			- Basiert auf NodeJS. NodeJs ist eine JavaScript Runtime, die JS Code serverseitig ausführen lässt.
			Das ist hier nicht weiter von relevanz.\\
			- Verwendet wird für die Anwendung ebenfalls roslibjs. Dies ist eine javascript bibliothek, welche in der html eingebunden wird. Sie stellt Befehle zum erstellen von ROS Subscriber/Publisher und Services bereit.\\
			
		\subsection{roslibjs}
		\label{Grundlagen:ROS2:RosLibJS}
	
	\section{DDS}
	\label{Grundlagen:DDS}
		- Grundlage für Kommunikation von ROS2\\
		- Realisiert die eigentliche Kommunikation\\
		- Wenn das hier geht, geht auch ROS2!
		- Verschiedene Systemanbieter, näher wird RTI und Fastrtps untersucht: https://ros.org/reps/rep-2000.html  (Beide TIER 1)
		
		ROS2 verwendet als Kommunikationsschnittstelle DDS.
		DDS funktioniert gut in lokalen Netzwerken, sobald allerdings ein NAT-Router dazwischen ist, funktioniert der standardmäßige Discovery Mechanismus nicht.
		Da das Ziel eine Edge Applikation ist, soll das Verhalten von DDS in der Docker Umgebung, bzw. in unterschiedlichen Netzwerken untersucht werden.
	
	\section{OpenController2}
	\label{Grundlagen:OpenController2}
		- Das ist ja meine zielhardware\\
		- Aktueller AGV läuft damit\\
		
	\section{TIA Openess}
	\label{Grundlagen:Fusionierung}
		- Lauras BA Vorstellen\\
		
		
		

	
	% \chapter{Voruntersuchungen}
\label{Voruntersuchungen}
	Im Rahmen dieser Masterarbeit werden mehrere Kommunikationswege über verschiedene Geräte verwendet.
	So ist es ein Ziel, das Framework ROS2 zu verwenden.
		
	\section{Testsetup}
	\label{Voruntersuchungen:Testsetup}
		Mithilfe einer Testumgebung können verschiedene Konfigurationen getestet und entwickelt werden.
		Notwendig ist, dass das Testsetup entscheidende Merkmale des Netzwerkes in der realen Umgebung beibehält.
		So ist die Aufteilung der Teilnehmer in zwei verschiedene Netzwerke wichtig, da die zu entwickelnde Edge App auf die Netzwerkeinstellung "host" verzichten soll.
		Ebenfalls wird dadurch eine sehr viel höhere Flexibilität der Netzwerkstruktur ermöglicht. 
				
		\subsection{Host Only}
		\label{Voruntersuchungen:Testsetup:HostOnly}
			Das Erste der beiden Testnetzwerke erstellt drei Virtuelle Maschinen: VM A, VM B und VM Router.
			Die Virtuelle Maschine VM A ist dem Subnetz 192.168.64.0/24 zugeordnet. VM B dem Subnetz 192.168.80.0/24.
			Als Brücke beider Subnetze dient VM Router, welche eine Verbindung zu beiden Netzwerken aufbaut und sämtliche Pakete an das jeweils andere Netz schickt.
			
			%\begin{table}[H]
			%	\centering
			%	\begin{tabular}{|l|p{2.5cm}|}
			%		\hline
			%		Virtuelle Maschine & IP Adresse\\
			%		\hline
			%		VM A & 192.168.64.20\\
			%		\hline
			%		VM B & 192.168.80.20\\
			%		\hline
			%		VM Router & 192.168.64.1  192.168.80.1 \\
			%		\hline					
			%	\end{tabular}
			%	\caption{\label{tab:VoruntersuchungenNetzwerke:Testsetup:HostOnly}IP Adressen des Host Only Setups.}
			%\end{table}
			\begin{figure}[H]
				\centering
				\input{"Bilder/tikz/Voruntersuchungen/struktur-host-only.tex"}
				\caption{Testsetup Host Only}
				\label{fig:Voruntersuchungen:Testsetup:HostOnly:Setup}					
			\end{figure}
			
			In Abbildung \ref{fig:Voruntersuchungen:Testsetup:HostOnly:Setup} sind die IP Adressen dieses Testsetups dargestellt.
					
		\subsection{Bridged}
		\label{Voruntersuchungen:Testsetup:Bridged}
		
			\begin{figure}[H]
				\centering
				\input{"Bilder/tikz/Voruntersuchungen/struktur-bridged.tex"}
				\caption{Testsetup Bridged}
				\label{fig:Voruntersuchungen:Testsetup:Bridged:Setup}					
			\end{figure}
			- Verwendet ein Bridged Netzwerk\\
			- Applikationen in Dockercontainern auf VM A\\
			- VM A und VM B im Bridged Netzwerk\
		
	\section{Lösungskonzepte}
	\label{Voruntersuchungen:Lösungskonzepte}
		\subsection{Initial Peers UDP}
		\label{Voruntersuchungen:Lösungskonzepte:InitialPeersUDP}
			- Im Discovery prozess von DDS wird eine Liste mit Partnern mitgegeben.\\
		
		\subsection{Initial Peers TCP}
		\label{Voruntersuchungen:Lösungskonzepte:InitialPeersTPC}
		 	- Eine liste mit TCP adressen $\rightarrow$ Voraussichtlich geht NAT einfacher\\
		
		\subsection{Server Client Discovery}
		\label{Voruntersuchungen:Lösungskonzepte:ServerClientDiscovery}
			- Ein eigener Discovery Server, der alle anmeldenden clients verbindet\\
		
		\subsection{Routing Service}
		\label{Voruntersuchungen:Lösungskonzepte:RoutingService}
			- Service von RTI, der gezielte Topics aus einer Domain in eine andere Domain überführen kann.\\
			- Andere Domain kann mittels TCP Stream WAN tauglich gemacht werden.\\
			- Quality of Service Einstellungen werden nicht automatisch erkannt. Bei besonderen Topics manuell einstellen. (Verweis auf meinen Forumsbeitrag ?).\\
			- Asymmetrische Konfiguration.\\
		

		
	\section{Erste Erkenntnisse}
	\label{VoruntersuchungenNetzwerke:Erkenntnisse}
		- Tabellarisch aufführen was geht und was nicht.
	
	\section{Proof of Concept}
	\label{VoruntersuchungenNetzwerke:ProofofConcept}
		Die in \ref{VoruntersuchungenNetzwerke:Erkenntnisse} vorgestellten Erkenntnisse sind für einen ersten Testlauf sehr interessant. Nachfolgend wird die Fähigkeit des Routing services anhand eines Proof of Concepts unter Beweis gestellt.
		Abbildung \ref{fig:VoruntersuchungenNetzwerke:ProofofConcept:Struktur} zeigt den prinzipiellen Versuchsaufbau.
		\begin{figure}[H]
			\centering
			\input{"Bilder/tikz/Voruntersuchungen/struktur-proof-of-concept.tex"}
			\caption{Setup für Proof of Concept}
			\label{fig:VoruntersuchungenNetzwerke:ProofofConcept:Struktur}					
		\end{figure}
		
		
		Auf dem Laptop sind mehrere Programme Installiert, so ist ein lokaler IEM eingerichtet, ein ROS2-Knoten, der die Tastatureingaben in ein \textit{/cmd\_vel} umwandelt.
		Anschließend wird das topic \textit{/cmd\_vel} in ein AGV spezifisches Topic der Form \textit{/cmd\_vel\_key/AGV*} umgewandelt, damit nur ein AGV gleichzeitig angesprochen wird.
		Diese Umwandlung geschieht in einen von Antonello Pastore bereitgestellten NodeRed Knoten.
		[SCREENSHOT]\\
		Die Routing Services von AGV 1 und AGV 2 bauen jeweils eine Verbindung zu dem Routing Service auf.
		Der Routing Service auf dem Laptop stellt für die AGV's die Topics \textit{*/cmd\_vel\_key/AGV*} bereit.
		Das '*' ist eine Wildcard und bildet die Filterregel, dass alle Topics, die \textit{/cmd\_vel\_key/AGV} haben, zu.
		Entsprechend erhält AGV 1 auch das nicht das AGV bestimmte Topic \textit{/cmd\_vel\_key/AGV2} und AGV 2 entsprechend \textit{/cmd\_vel\_key/AGV1}.
		Ein ROS2 Demultiplexer, der AGV spezifisch eingestellt ist, filtert nun das für das AGV relevante topic und setzt das topic \textit{/cmd\_vel\_key/AGV1} auf \textit{/cmd\_vel\_key} um.
		In den Eingang des twist multiplexers, dessen Ausgang die Bewegungssteuerung des AGV's vorgibt, ist nun \textit{/cmd\_vel\_key} als zusätzlicher Eingang konfiguriert.
		
		Dieses Setup ermöglicht eine Auswahl des AGV's, welcher vom Laptoop aus gesteuert werden soll und behält die Steuerbarkeit mittels Joystick, welcher als Bluetooth Gerät im internen Netzwerk des AGV's angeschlossen ist bei.

		
	
		
		
	
	
	% \chapter{Konzeption und Implementierung}
\label{Implementierung}
	\section{Struktur}
	\label{Implementierung:Struktur}
	- Präsentation, die ich laura gezeigt habe\\
	\begin{figure}[H]
		\centering
		\input{"Bilder/tikz/Implementierung/struktur-gesamt.tex"}
		\caption{Überblick über die Struktur der Implementierung}
		\label{fig:Implementierung:Struktur:Gesamt}					
	\end{figure}
	Hier verwende ich mal eine Referenz zu \cite{Krypczyk2018}
	\begin{figure}[H]
		\centering
		\input{"Bilder/tikz/Implementierung/Sequenzdiagram.tex"}
		\caption{Test des latex-Paketes pgf-umlsd}
		\label{fig:Implementierung:Struktur:pgf-umlsd-test}		
	\end{figure}
	asdf
	\begin{figure}[H]
		\centering
		\input{"Bilder/tikz/Implementierung/GesamtDiagramm.tex"}
		\caption{Überblick über die Struktur der Implementierung}
		\label{fig:Implementierung:Struktur:GesamtDiagramm}					
	\end{figure}
	\section{Device}
	Hallo, ich bin \faCcVisa
	\label{Implementierung:Device}
	
		\subsection{Struktur}
		\label{Implementierung:Device:Struktur}
	
		\subsection{Masterapp}
		\label{Implementierung:Device:Masterapp}
		
		\subsection{Clientapp}
		\label{Implementierung:Device:Clientapp}
		
	\section{Server}
	\label{Implementierung:Server}
		
		\subsection{Anpassungen an Fusionsprojekt}
		\label{Implementierung:Server:Anpassungen}
		
		\subsection{Struktur}
		\label{Implementierung:Server:Struktur}
		
		\subsection{Tbd}
		
	\section{Userinterface}
	\label{Implementierung:Userinterface}
	
		\subsection{Design}
		\label{Implementierung:Userinterface:Design}
	
		\subsection{Kommunikationsschnittstellen}
		\label{Implementierung:Userinterface:Kommunikationsschnittstellen}
		
		
		
	
	
	
	% \chapter{Proof of Concept}
\label{PoC}

	\section{Manueller Test}
	\label{Design:Struktur}
		- Prozess soll von Laien ausführbar sein
	
	\section{Integration in Edge}
	\label{Design:Abhängigkeiten}
	
	\section{Fusionierung mehrerer Projekte}
	\label{PoC:Fusionierung}	
	
	% \chapter{Ausblick}
\label{Ausblick}
	- Schnittstelle für TIA Projekte definieren\\
	- Funtionserweiterung für TIA Openess\\
	- Device seite unabhängig von Master Client\\
	
	
	
	% \chapter{Fazit}
\label{Fazit}
	Mal schauen was so kommt :-)
	
	\label{Ende} % Wichtig für Fußnote "Seite X von Y"	
	\clearpage
			
	% Für Anhang wieder umstellen auf Römische Seitenzähler

	\pagenumbering{Roman}
	\setcounter{page}{\value{SeitenzahlSpeicher}}
	\ofoot{\pagemark}
	
	% Literaturverzeichnis
	\printbibliography[title={Literaturverzeichnis}]
	
	% Abbildungsverzeichnis
	\listoffigures 
	
	% Tabellenverzeichnis
	\listoftables
	
	% Listing-Verzeichnis
	% Namen definieren
	% \renewcommand{\lstlistlistingname}{Verzeichnis der Listings}
	\lstlistoflistings
	
	
\end{document}